\documentclass[11pt,a4paper]{article}
\usepackage{genesisl1-publication}
\GenesisTitle{Title of the Community Letter}
\GenesisShortTitle{Community letter}
\GenesisSubtitle{Optional human-centred subtitle}
\GenesisType{Community Letter}
\GenesisAuthor{Author or contributor group}
\GenesisAffiliation{Role, community, or independent contributor}
\GenesisVersion{Draft 0.1}
\GenesisStatus{OPEN FOR DISCUSSION}
\GenesisIdentifier{Optional document ID}
\GenesisLicense{CC BY 4.0}
\begin{document}
\MakeGenesisTitle
\begin{GenesisAbstract}
Explain why you are writing and what kind of response, reflection, or participation would
be helpful.
\end{GenesisAbstract}
\section{Dear community}
Open in your own voice. A community letter may be personal, reflective, technical,
philosophical, celebratory, corrective, or invitational.
\section{Why I am writing}
Explain the moment, concern, opportunity, or appreciation that motivated the letter.
\section{What I have heard and learned}
Acknowledge contributions and perspectives fairly. This optional section can show how
community discussion influenced the message.
\section{Present understanding}
Describe your current view without presenting it as the only acceptable position.
\ClaimStatus{Status}{Reflection, proposal, interpretation, update, or decision. Label it honestly.}
\section{What remains open}
Identify questions, disagreements, uncertainties, and decisions that still need collective work.
\section{Proposed direction}
Offer a direction without forcing contributors into one writing style, discipline, governance
model, or technical framework.
\begin{GenesisAction}[Invitation to the community]
Request the forms of participation that would be useful: critique, references, experiments,
design alternatives, implementation, translation, documentation, or reflection. Welcome
disagreement and preserve contributors' dignity.
\end{GenesisAction}
\section*{Acknowledgements}
Thank contributors with their consent.
\end{document}
